\documentclass[12pt]{article}
\usepackage[utf8]{inputenc}
\usepackage[T1]{fontenc}
\usepackage{hyperref}
\title{Hdi Language Resources for K-12 Curriculum Development}
\author{Compiled by Claude Sonnet 5.5 (claude-sonnet-5-5) for Andr\'as Kornai}
\date{}
\begin{document}
\maketitle

\section{Introduction}
Hdi (ISO 639-3 xed; Glottocode hdii1240) is a Central Chadic language of northern Cameroon and northeastern Nigeria. Most speakers are in Nigeria. In Cameroon the community is concentrated in Tourou, near the Nigerian border. It is also called Hide, Hedi, Turu-Hide or Tur, and speakers call it xd\'i. Speaker estimates vary: Wikipedia cites about 29,000 from 2001, while the dictionary's authors give 10,000 to 25,000. It is written in the Latin alphabet. I found no school materials for Hdi.

\section{Primary Educational Resources}
None found in the search results. The authors of the 2015 dictionary hoped it would help speakers learn English and French.

\section{Linguistic and Reference Materials}
\begin{itemize}
\item \emph{Dictionary of Hdi} (2015), edited by Zygmunt Frajzyngier and colleagues. It is trilingual (Hdi--English--French) and includes a grammatical sketch. Publisher page: \url{https://www.koeppe.de/titel_dictionary-of-hdi}
\item Print edition: \url{https://www.koeppe.de/titel_print_dictionary-of-hdi}
\item Hdi dictionary on Dictionaria: \url{https://dictionaria.clld.org/contributions/hdi}
\item News item on the dictionary, University of Colorado Boulder: \url{https://www.colorado.edu/linguistics/2016/01/25/zygmunt-frajzyngier-publishes-new-dictionary-chadic-language}
\item Hdi language (Wikipedia): \url{https://en.wikipedia.org/wiki/Hdi_language}
\item The Hide or Tur: \url{https://mandaras.info/HideTur.html}
\end{itemize}

\section{Religious and Community Resources}
None found in the search results.

\section{Digital and Audio Resources}
\begin{itemize}
\item Omniglot page on Hdi writing: \url{https://omniglot.com/writing/hdi.htm}
\end{itemize}
The dictionary is reported to work with the Lexus database, which has sound recordings. I did not see a URL for that database.

\section{Curriculum Development Context}
The 2015 trilingual dictionary and its grammatical sketch are the main base for any curriculum. The search results show no existing school instruction, primers or community education organization, so these would need to be developed with the community.

\end{document}